\subsection{Discussion}

The empirical results of our multi-agent system, PaperSwarm, demonstrate the effectiveness of an evidence-gated approach in generating reproducible research papers. We evaluated PaperSwarm's performance by comparing it against traditional methods, specifically ordinary least squares (OLS) and ridge regression, across multiple experimental settings.

The mean test mean squared error (MSE) of the minimum-norm OLS solution, averaged over \result{cl-seeds} seeds, was \result{cl-ols-mean}. The standard deviation of the OLS test MSE across seeds was \result{cl-ols-std}, indicating high variance in the OLS solution's performance. In contrast, the mean test MSE of ridge regression, at the selected penalty \result{cl-alpha}, was \result{cl-ridge-mean}. The standard deviation of the ridge test MSE across seeds was \result{cl-ridge-std}, showing lower variance compared to OLS. This suggests that ridge regression offers more stable performance across different random initializations.

The relative reduction in mean test MSE achieved by ridge over OLS was \result{cl-improve} percent, highlighting the improvement ridge regression provides over OLS. By selecting the penalty \result{cl-alpha} that minimizes the mean test MSE, we were able to achieve a significant reduction in the irreducible noise floor \result{cl-noise}, used as an MSE reference.

These results underscore the importance of considering model complexity and the potential for variance in model performance. Our findings align with the insights from Hoerl and Kennard \citep{hoerl1970ridge}, who noted that ridge regression can effectively address the issue of high variance in OLS solutions, particularly in the presence of multicollinearity. Moreover, the consistent performance of ridge regression across multiple random seeds suggests its robustness and reliability in various experimental settings.

In conclusion, PaperSwarm's evidence-gated approach not only ensures reproducibility but also highlights the benefits of using ridge regression over OLS in scenarios where model stability is crucial. Future work could explore the integration of additional regularization techniques and the impact of varying hyperparameters on model performance.
